% Part: normal-modal-logic
% Chapter: syntax-and-semantics
% Section: entailment

\documentclass[../../../include/open-logic-section]{subfiles}

\begin{document}

\olfileid[hi]{nml}{syn}{ent}

\olsection{परिणाम}

\begin{explain}
  किसी जगत पर सत्यता की परिभाषा से हम !!{formula}ों के बीच परिणाम संबंध
  परिभाषित कर सकते हैं। कोई !!{formula}~$!B$, $!A$ को तभी परिणामतः देता है
  जब $!B$ के सत्य होने पर $!A$ भी सत्य हो। यहाँ ``जब भी'' का अर्थ है
  ``हम चाहे जिस प्रतिरूप पर विचार करें'' और ``उस प्रतिरूप के चाहे जिस जगत
  पर विचार करें'' दोनों।
\end{explain}

\begin{defn}
  यदि $\Gamma$ !!{formula}ों का समुच्चय और $!A$ कोई !!{formula} है, तो
  $\Gamma$ से~$!A$ \emph{परिणामतः मिलता है}, संकेत में
  $\Gamma \Entails !A$, तभी जब प्रत्येक प्रतिरूप
  $\mModel{M} = \tuple{W, R, V}$ और जगत $w \in W$ के लिए, यदि प्रत्येक
  $!B \in \Gamma$ के लिए $\mSat{M}{!B}[w]$, तो $\mSat{M}{!A}[w]$। यदि
  $\Gamma$ में केवल एक !!{formula}~$!B$ है, तो हम $!B \Entails !A$ लिखते हैं।
\end{defn}

\begin{ex}
  यह दिखाने के लिए कि किसी !!{formula} से दूसरा परिणामतः मिलता है, हमें
  $\mSat{M}{}[w]$ की परिभाषा का उपयोग करके सभी प्रतिरूपों के विषय में तर्क
  करना होता है। उदाहरणार्थ, $p \lif \Diamond p \Entails
  \Box\lnot p \lif \lnot p$ दिखाने के लिए हम इस प्रकार तर्क कर सकते हैं:
  कोई प्रतिरूप $\mModel{M} = \tuple{W, R, V}$ और $w \in W$ लें और मान लें
  $\mSat{M}{p \lif \Diamond p}[w]$। हमें
  $\mSat{M}{\Box\lnot p \lif \lnot p}[w]$ दिखाना है। मान लें ऐसा नहीं है।
  तब $\mSat{M}{\Box\lnot p}[w]$ और $\mSat/{M}{\lnot p}[w]$। चूँकि
  $\mSat/{M}{\lnot p}[w]$, इसलिए $\mSat{M}{p}[w]$। धारणा से
  $\mSat{M}{p \lif \Diamond p}[w]$, अतः $\mSat{M}{\Diamond p}[w]$।
  $\mSat{M}{\Diamond p}[w]$ की परिभाषा से कोई $Rww'$ वाला $w'$ ऐसा है कि
  $\mSat{M}{p}[w']$। चूँकि $\mSat{M}{\Box \lnot p}[w]$ भी है, इसलिए
  $\mSat{M}{\lnot p}[w']$, जो विरोधाभास है।

  यह दिखाने के लिए कि किसी !!{formula}~$!B$ से दूसरा~$!A$ परिणामतः नहीं
  मिलता, हमें प्रतिदृष्टांत देना होता है: ऐसा प्रतिरूप
  $\mModel{M} = \tuple{W, R, V}$ जिसमें किसी जगत~$w \in W$ पर
  $\mSat{M}{!B}[w]$ पर $\mSat/{M}{!A}[w]$ हो। आइए दिखाएँ कि
  $p \lif \Diamond p \Entails/ \Box p \lif p$।
  \olref{fig:counterex} के प्रतिरूप पर विचार करें। हमारे पास
  $\mSat{M}{\Diamond p}[w_1]$, अतः
  $\mSat{M}{p \lif \Diamond p}[w_1]$। किंतु
  $\mSat{M}{\Box p}[w_1]$ पर $\mSat/{M}{p}[w_1]$ होने के कारण
  $\mSat/{M}{\Box p \lif p}[w_1]$।
  \begin{figure}
  \begin{center}
    \begin{tikzpicture}[modal]
      \node[world] (w1) [label=right:\mFalse{p}]{$w_1$}; 
      \node[world] (w2) [label=right:\mTrue{p}, above left=of w1]{$w_2$}; 
      \node[world] (w3) [label=right:\mTrue{p}, above right=of w1] {$w_3$};
      \draw[->] (w1) to (w2);
      \draw[->] (w1) to (w3);
    \end{tikzpicture}
  \end{center}
  \caption{$p \lif \Diamond p \Entails \Box p \lif p$ का प्रतिदृष्टांत।}
  \ollabel{fig:counterex}
  \end{figure}
  
  बहुत सरल प्रतिदृष्टांत प्रायः पर्याप्त होते हैं। $W' = \{w\}$,
  $R' = \emptyset$ और $V'(p) = \emptyset$ वाला प्रतिरूप
  $\mModel{M'} = \{W', R', V'\}$ भी प्रतिदृष्टांत है: चूँकि
  $\mSat/{M'}{p}[w]$, इसलिए $\mSat{M'}{p \lif \Diamond p}[w]$। $w$ से
  कोई जगत अभिगम्य न होने के कारण $\mSat{M'}{\Box p}[w]$, अतः
  $\mSat/{M'}{\Box p \lif p}[w]$।
\end{ex}

\begin{prob}
  दिखाइए कि $\Box (!A \land !B) \Entails \Box !A$।
\end{prob}

\begin{prob}
  दिखाइए कि $\Box (p \lif q) \Entails/ p \lif \Box q$ और
  $p \lif \Box q \Entails/\Box (p \lif q)$।
\end{prob}

\end{document}
